\section{Flow model and the fixed approximation space}
\label{sec:governing-equations}
\label{sec:space-discretization}

\subsection{A driven finite-annulus problem}

We consider the two-dimensional incompressible Navier--Stokes equations
\begin{equation}
 \partial_t\boldsymbol u+(\boldsymbol u\cdot\nabla)\boldsymbol u
 =-\nabla p+\nu\Delta\boldsymbol u,\qquad
 \nabla\cdot\boldsymbol u=0,\qquad
 \Omega=\{(r,\theta):a<r<R\}.
 \label{eq:ns}
\end{equation}
The cylinder diameter and the final streamwise speed define the length
and velocity scales. The calculations use \(a=0.5\), \(R=6\),
\(\nu=0.01\), and hence \(\mathrm{Re}=100\). The final time is
\(T=72\). Introduce the logarithmic radial coordinate
\(s=\log(r/a)\), its extent \(L=\log(R/a)\), and \(z=s/L\).
The area measure becomes \(r^2\,ds\,d\theta\). The physical radial span is $R-a=5.5$; it is not the logarithmic extent $L=\log12$ used in the formulas.

For nonzero angular modes, a streamfunction gives
\begin{equation}
 u_r=\frac{\psi_\theta}{r},\qquad
 u_\theta=-\frac{\psi_s}{r},\qquad
 \omega=-\frac{\psi_{ss}+\psi_{\theta\theta}}{r^2}.
 \label{eq:kinematics}
\end{equation}
The azimuthal mean is advanced separately as a real tangential velocity
\(v(s,t)\); its radial velocity is zero and its vorticity is
\((v+v_s)/r\). This representation preserves all sine and cosine
components of every retained nonzero mode and does not impose a
reflection or zero-circulation constraint. The pressure is eliminated
from the nonzero-mode equations by taking the curl; the mean tangential
momentum equation supplies the remaining velocity block.

The initially stationary flow is driven by
\begin{equation}
 U(t)=1-e^{-(t/0.8)^3},\qquad
 V(t)=\begin{cases}
 0.12\sin^4[\pi(t-4)/8],&4<t<12,\\
 0,&\text{otherwise}.
 \end{cases}
 \label{eq:driving}
\end{equation}
The finite-duration transverse pulse breaks reflection symmetry. At
the inner boundary, \(\psi=\psi_s=0\). At the outer boundary, the
nonzero modes satisfy
\begin{equation}
 \psi(L,\theta,t)=R[U(t)\sin\theta-V(t)\cos\theta],\qquad
 \omega(L,\theta,t)=0.
 \label{eq:outer-boundary}
\end{equation}
The separate mean block satisfies \(v(0,t)=0\) and
\(v_s(L,t)+v(L,t)=0\). Thus the outer circle does not have a full
velocity Dirichlet condition. It is part of the definition of the
finite-domain problem, rather than an exact transparent boundary for
unbounded cylinder flow.

The prescribed boundary data are represented by a lifting
\begin{equation}
 \begin{gathered}
 \psi_\ell(s,\theta,t)=\ell(s)[U(t)\sin\theta-V(t)\cos\theta],\\
 \ell(s)=R[(1-\beta)z^2+\beta z^3],\quad
 \beta=\frac{L^2-2}{4}.
 \end{gathered}
 \label{eq:lifting}
\end{equation}
It obeys \(\ell(0)=\ell_s(0)=0\), \(\ell(L)=R\), and
\(\ell_{ss}(L)=R\). The last identity enforces zero outer vorticity
for the lifted \(k=1\) mode. With the convention
\(f=c_0+2\operatorname{Re}\sum_{k=1}^{K}c_ke^{ik\theta}\),
the complex lifting amplitude is \(g(t)=-(V(t)+iU(t))/2\).
The same convention is used in reconstruction, nonlinear projection,
and error scoring; conjugate negative modes represent the complete
real field.

The finite outer radius is a material limitation. In the archived
domain-size diagnostic, increasing \(R\) from 6 to 8 changed the
unaligned vorticity by about \(50.39\%\) in the common \(r\leq4\)
region, and also changed the shedding frequency. The numerical
tolerance below is therefore a tolerance for the specified
\(R=6\) problem. Domain truncation and discretization error are not
interchangeable.

\subsection{Random and deterministic radial trial functions}

Write the nonzero streamfunction coefficients as
\begin{equation}
 \psi_k(s,t)=\sum_{j=1}^{d}F_j(s)x_{jk}(t)
                +\delta_{k1}g(t)\ell(s),\qquad 1\leq k\leq K.
 \label{eq:trial-expansion}
\end{equation}
The functions \(F_j\) are fixed before time integration. No neural
network weights are trained along the trajectory. The RFM designation
refers to the randomly chosen radial feature space, followed by a
deterministic Galerkin evolution of its coefficients \cite{r1,r2,r3}.

For an odd input size \(m=2p+1\), the raw random features are a constant
and \(p\) sine--cosine pairs. Specifically,
\begin{equation}
 \begin{gathered}
 b_j=\frac{62}{p}(j+\xi_j),\qquad
 \xi_j\sim\operatorname{Uniform}(0.2,0.8),\\
 q_{j,0}(z)=\cos[b_j(z-\tfrac12)],\quad
 q_{j,1}(z)=\sin[b_j(z-\tfrac12)].
 \end{gathered}
 \label{eq:random-features}
\end{equation}
where \(j=0,\ldots,p-1\). The random draws use a fixed
\texttt{NumPy default\_rng(seed)}. The registered seed set is
\(7331,7687,8011,8353,9041\). The Legendre comparison replaces these
raw functions by \(P_j(2z-1)\), \(j=0,\ldots,m-1\), while retaining
the same boundary transformation, enrichment, quadrature, and time
discretization. The reported cases use no radial coordinate warping.

Boundary conditions are imposed analytically. For a raw function
\(q\), let \(\chi(z)=16z^2(1-z)\) and
\(b(z)=z^2(1-z)^2/2\), and define
\begin{equation}
 \mathcal Bq(z)=\chi(z)q(z)
       -b(z)\left.\frac{d^2(\chi q)}{dz^2}\right|_{z=1}.
 \label{eq:boundary-transform}
\end{equation}
Because \(b(0)=b'(0)=b(1)=0\) and \(b''(1)=1\), this function
satisfies \(F(0)=F_s(0)=F(L)=F_{ss}(L)=0\). Consequently the
homogeneous nonzero modes satisfy both no slip at the cylinder and
zero outer vorticity for every \(k\). Derivatives with respect to
\(s\) include the factors \(L^{-1}\) and \(L^{-2}\); the
correction in \eqref{eq:boundary-transform} is expressed in \(z\).

Both trial families receive the same six analytic outer-layer
functions. For \(\delta=L-s\), these are
\begin{equation}
 E_\gamma(s)=z^2\left[\delta+
       \left(\gamma+\frac{2}{L}\right)\delta^2\right]
       e^{-\gamma\delta},\qquad
 \gamma\in\{16,36,80,165,340,660\}.
 \label{eq:layer-features}
\end{equation}
Their constant and quadratic Taylor coefficients at \(\delta=0\)
vanish, giving the same outer homogeneous conditions; the factor
\(z^2\) supplies the inner conditions. These functions resolve the
outer-boundary layers of this finite-domain formulation. They are not
a learned consequence of the random features.

The mean block uses a deterministic polynomial space in both families.
Its raw functions are
\(z[(L+2)-(L+1)z]\) and
\(z(1-z)^2P_j(2z-1)\), \(j=0,\ldots,m_0-2\), with \(m_0=24\).
The mean enrichment is
\(z[1+(\gamma+1+1/L)\delta]e^{-\gamma\delta}\).
All satisfy the mean boundary conditions. The six enrichments give
30 retained mean functions in the final configurations. Randomness is
therefore restricted to the nonzero-mode radial trial space.

\subsection{Quadrature, rank selection, and the meaning of model size}

Radial integration uses positive composite Gauss weights on
\([0,L-0.12]\), \([L-0.12,L-0.018]\), and \([L-0.018,L]\).
With \(q\) nodes and \(b_q=\max(8,\lfloor q/6\rfloor)\), the
three allocations are \(q-2b_q,b_q,b_q\). For the two principal
Galerkin configurations, these are \(44/11/11\) and \(52/12/12\).
The random trigonometric and exponential radial functions are not
polynomials, so this quadrature is an approximation to the continuous
weak form. Positivity statements below apply exactly to the
quadrature-defined matrices; they are not quadrature-error estimates.

Let \(\widetilde F,\widetilde D,\widetilde D_2\) be the values and
first two \(s\)-derivatives of all raw, boundary-corrected and enriched
functions at the radial nodes. With \(W=\operatorname{diag}(w_i)\),
form
\begin{equation}
 \mathcal A=\operatorname{diag}(\sqrt{w_i}/r_i)
              (\widetilde D_2-\widetilde F),\quad
 S=\operatorname{diag}(\|\mathcal A_{:j}\|_2),\quad
 \mathcal A S^{-1}=U\Sigma V^{\mathsf T}.
 \label{eq:feature-svd}
\end{equation}
All retained raw columns must have nonzero norms. Keep singular values
\(\sigma_j>10^{-10}\sigma_1\), and set
\begin{equation}
 C=S^{-1}V_d\Sigma_d^{-1},\qquad
 F=\widetilde F C,\quad D=\widetilde D C,\quad
 D_2=\widetilde D_2 C.
 \label{eq:feature-transform}
\end{equation}
This is a discrete differential normalization with rank truncation,
rather than an orthogonalization in the ordinary point-value norm.
In exact arithmetic,
\(\operatorname{diag}(\sqrt{w_i}/r_i)(D_2-F)=U_d\).
It therefore normalizes the \(k=1\) dissipation Gram matrix to
\(\nu I\). This identity describes the existing construction; it is
not a newly measured optimization or a guarantee for all angular modes.
The mean block analogously normalizes
\(\operatorname{diag}(\sqrt{w_i})(\widetilde V+\widetilde V_s)\).

An input label such as RFM37 counts raw random features, not unknowns.
Enrichment adds columns and truncation can remove them. The actual
nonzero-mode rank is \(d\), and the real coefficient count is
\(2Kd+d_0\), where \(d_0=30\) is the mean rank. Fixed-rank contrasts
and accuracy-matched contrasts answer different questions. In
particular, a lower RFM input count is not evidence of a lower rank
than a Legendre control. The seed and cutoff are part of the model
identity and must remain fixed when comparing implementations.

\section{Weak formulation and time-discrete algorithms}
\label{sec:numerical-method}

\subsection{Nonzero-mode mass, dissipation, and nonlinear flux}

For the velocity field defined above, incompressibility gives the
vorticity equation in conservative logarithmic coordinates:
\begin{equation}
 r^2\partial_t\omega
  +\partial_s(r u_r\omega)+\partial_\theta(r u_\theta\omega)
   =\nu(\omega_{ss}+\omega_{\theta\theta}).
 \label{eq:vorticity-conservative}
\end{equation}
Test its \(k\)-th Fourier coefficient with a homogeneous radial basis
function. One integration by parts in the mass term after inserting
\(r^2\omega_k=-(\psi_{k,ss}-k^2\psi_k)\) yields the bilinear form
\(\int(F_i'\psi_k'+k^2F_i\psi_k)\,ds\). Integrating the viscous
term twice gives a negative Gram form. The boundary terms vanish
because \(F_i=0\) at both ends, \(F_i'(0)=0\), and
\(\omega_k(L)=0\). Applying the prescribed radial quadrature gives
\begin{align}
 M_k&=D^{\mathsf T}WD+k^2F^{\mathsf T}WF,
       \label{eq:mass}\\
 K_k&=\nu(D_2-k^2F)^{\mathsf T}
        \operatorname{diag}(w_i/r_i^2)(D_2-k^2F).
       \label{eq:dissipation}
\end{align}
For every real coefficient vector \(x\),
\begin{equation}
 \begin{aligned}
 x^{\mathsf T}M_kx&=\|W^{1/2}Dx\|_2^2+k^2\|W^{1/2}Fx\|_2^2,\\
 x^{\mathsf T}K_kx&=\nu\|
    \operatorname{diag}(\sqrt{w_i}/r_i)(D_2-k^2F)x\|_2^2.
 \end{aligned}
 \label{eq:gram-positivity}
\end{equation}
Both matrices are symmetric positive semidefinite in exact
arithmetic. Positive definiteness of \(M_k\) additionally requires
full column rank of the stacked weighted value and derivative samples
on the retained space. The compact form \eqref{eq:dissipation} is
algebraically equivalent to the expanded expression used by the
archived implementation. Its positivity does not remove finite
precision cancellation, assembly asymmetry, or a poorly conditioned
coordinate map.

Let \(a_k\) and \(b_k\) be the Fourier coefficients of
\(u_r\omega\) and \(u_\theta\omega\) evaluated at the radial
nodes. Integration by parts in the two transport terms gives
\begin{equation}
 f_k=D^{\mathsf T}\operatorname{diag}(w_i r_i)a_k
      -ikF^{\mathsf T}\operatorname{diag}(w_i r_i)b_k.
 \label{eq:weak-flux}
\end{equation}
The products use the complete velocity, including the mean block and
boundary lifting. They are evaluated in physical angle and projected
back to the retained spectrum, coupling the otherwise independent
linear mode blocks. This is the reason that making linear solves
diagonal does not eliminate the nonlinear FFT work.

For \(k=1\), the lifting projections are
\begin{equation}
 \begin{aligned}
 M_\ell&=D^{\mathsf T}W\ell_s+F^{\mathsf T}W\ell,\\
 K_\ell&=\nu(D_2-F)^{\mathsf T}
       \operatorname{diag}(w_i/r_i^2)(\ell_{ss}-\ell).
 \end{aligned}
 \label{eq:lift-projections}
\end{equation}
They vanish from the equations of the other modes. Suppressing the
mode index and allowing zero lifting vectors where appropriate, the
semidiscrete system is
\begin{equation}
 m'(t)+Kx(t)+g(t)K_\ell=f(t),\qquad
 m(t)=Mx(t)+g(t)M_\ell.
 \label{eq:semidiscrete}
\end{equation}
The differentiated quantity is the complete weak mass \(m\). Treating
it as \(Mx\) alone would omit the boundary acceleration. Several
independent liftings can be included by replacing each \(gM_\ell\)
and \(gK_\ell\) by the corresponding sum.

\subsection{Why the mean block is a separate equation}

Angular averaging of tangential momentum gives
\begin{equation}
 \partial_t v=-\langle u_r\omega\rangle_\theta
       +\frac{\nu}{r^2}(v_{ss}-v).
 \label{eq:mean-equation}
\end{equation}
Let \(V\) and \(V_s\) now denote the sampled mean basis and its
derivative, rather than the transverse pulse in
\eqref{eq:driving}. The mean matrices and source are
\begin{equation}
 \begin{aligned}
 M_0&=V^{\mathsf T}\operatorname{diag}(w_i r_i^2)V,\qquad
 K_0=\nu(V+V_s)^{\mathsf T}W(V+V_s),\\
 f_0&=-V^{\mathsf T}\operatorname{diag}(w_i r_i^2)a_0.
 \end{aligned}
 \label{eq:mean-matrices}
\end{equation}
To check the dissipation form, expand
\(\int(V+V_s)(v+v_s)\,ds\): its cross terms combine to
\([Vv]_0^L\), which is exactly the boundary contribution obtained
from \(-\int V(v_{ss}-v)\,ds\) using \(v_s(L)=-v(L)\) and
\(v(0)=0\). There is no direct boundary lifting in this block.
The negative radial-flux sign in \eqref{eq:mean-matrices} follows
from tangential momentum. Setting \(k=0\) in
\eqref{eq:mass}--\eqref{eq:weak-flux} would produce a different
representation and is not the implemented mean equation.

\subsection{Constant-step integration and a single bridge}
\label{sec:time-phase}

The baseline treats diffusion implicitly and the nonlinear weak flux
by two-point extrapolation, using SBDF2 \cite{r9,r12}. After the
implicit Euler start, a constant step \(h\) solves
\begin{equation}
 \begin{aligned}
 (\tfrac32M+hK)x_{n+1}
  ={}&2m_n-\tfrac12m_{n-1}+h(2f_n-f_{n-1})\\
   &-g_{n+1}(\tfrac32M_\ell+hK_\ell).
 \end{aligned}
 \label{eq:uniform-sbdf}
\end{equation}
The initialization uses the stationary initial state and the
prescribed smooth startup. Euler startup does not establish
second-order convergence by itself; its consistency with the
initial-value problem must be considered together with the following
multistep evolution. In the implemented initial state the explicit
nonlinear source is zero.

The optimized schedule uses a constant step \(H\) until \(t_0=8\)
and a larger constant step \(h\) afterward. Only the first update
across this interface has unequal consecutive intervals. Write
\(t_-=t_0-H\), \(t_+=t_0+h\), and \(\rho=h/H>0\). Define
\begin{equation}
 a_0=\frac{1+2\rho}{1+\rho},\qquad b_0=1+\rho,\qquad
 b_1=\frac{\rho^2}{1+\rho}.
 \label{eq:bridge-coefficients}
\end{equation}
The bridge is
\begin{equation}
 \begin{aligned}
 (a_0M+hK)x_+={}&b_0m_0-b_1m_-+
 h[(1+\rho)f_0-\rho f_-]\\
 &-g_+(a_0M_\ell+hK_\ell).
 \end{aligned}
 \label{eq:bridge}
\end{equation}
In this subsection \(m_0\) and \(f_0\) refer to time \(t_0\),
not to the azimuthal mean. Historical masses already contain their
historical lifting values; the future lifting is subtracted once.
At \(\rho=1\), \eqref{eq:bridge} reduces exactly to
\eqref{eq:uniform-sbdf}.

\paragraph{Proposition 1 (local consistency of the lifted bridge).}
Assume \(m\in C^4\) and \(f\in C^3\) along a smooth exact
semidiscrete solution in a neighborhood of the bridge. Assume also
that \(\rho\) lies in a fixed compact subset of \((0,\infty)\).
For \(\delta=\max(H,h)\), the derivative-form residual of
\eqref{eq:bridge} is \(O(\delta^2)\), and its step-scaled
equation residual is \(O(h\delta^2)\).

\paragraph{Proof.}
Differentiate the quadratic interpolant through the three masses at
\(t_+\). Its derivative is
\(Dm_+=(a_0m_+-b_0m_0+b_1m_-)/h\).
The linear extrapolant of the two fluxes has value
\((1+\rho)f_0-\rho f_-\) at the same endpoint. Taylor expansion
about \(t_+\) gives
\begin{align}
 Dm_+-m'(t_+)
   &=-\frac{h(h+H)}{6}m'''(t_+)+O(\delta^3),
      \label{eq:bridge-derivative}\\
 (1+\rho)f_0-\rho f_- - f(t_+)
   &=-\frac{h(h+H)}{2}f''(t_+)+O(\delta^3).
      \label{eq:bridge-extrapolation}
\end{align}
Substitute these in \eqref{eq:semidiscrete}, evaluate diffusion and
lifting at \(t_+\), and collect the two residuals. Substitution of
\(m_+=Mx_++g_+M_\ell\) gives the stated bridge. \hfill\(\square\)

This proposition is local consistency, not unconditional stability
or long-time convergence of the nonlinear Navier--Stokes solution.
Those require additional regularity, startup, perturbation, and
stability estimates. The finite-duration pulse is piecewise smooth;
the displayed Taylor argument is local to the bridge at \(t=8\),
which is away from the pulse endpoints. For the scalar zero equation
\(y'=0\), the bridge differences satisfy
\begin{equation}
 y_+-y_0=\frac{\rho^2}{1+2\rho}(y_0-y_-).
 \label{eq:zero-equation}
\end{equation}
The two registered ratios \(8/7\) and \(19/17\) give
\(64/161\) and \(361/935\); the constant segments have factor
\(1/3\). Thus the homogeneous differences are bounded on this
particular grid. Explicit advection and mode coupling are absent
from this check, so it does not supply a flow stability theorem.

The work reduction follows from the schedule rather than from the
trial family. At \(T=72\), changing the rate from 128 to 112 at
\(t=8\) changes the step count from \(9216\) to \(8192\).
Changing 152 to 136 changes \(10944\) to \(9920\). The initial
nonlinear source is identically zero and the Euler state is constructed
before the repeated updates. Nonlinear kernels evaluate the states
at \(t_1,\ldots,t_{N_{\rm step}-1}\); the terminal state needs no
subsequent evaluation. The respective
recorded evaluation counts are \(9215\to8191\) and
\(10943\to9919\). This count does not predict the complete solve
time, which also includes additional phase preparation and startup.

\subsection{Reuse of fixed spatial preparation}

\paragraph{Proposition 2 (exact-arithmetic reuse).}
Fix the geometry, viscosity, feature family, random seed, enrichments,
quadrature, lifting, column scaling, and singular-value cutoff. If the
basis transformation is formed before any step-dependent matrix and
does not depend on the time step, then the retained bases, ranks,
\(M,K,M_\ell,K_\ell\), and uncontracted nonlinear test operators
can be prepared once for all time phases without changing the
discrete update in exact arithmetic.

\paragraph{Proof.}
Equations \eqref{eq:feature-svd}--\eqref{eq:feature-transform} and
\eqref{eq:mass}--\eqref{eq:mean-matrices} contain no time-step
parameter. Phase dependence enters only through \(a_jM+h_jK\),
its application or factorization, and the step-weighted right-hand
side. Reusing the same inputs to these operations therefore gives
the same linear equations. \hfill\(\square\)

The existing assembly path accepts
\(A_{\rm eff}=(3/2)M+h_{\rm eff}K\). For the bridge set
\(c=(3/2)/a_0\) and \(h_{\rm eff}=ch\). Then
\begin{equation}
 A_{\rm eff}=c(a_0M+hK),\qquad
 A_{\rm eff}^{-1}(c b)=(a_0M+hK)^{-1}b.
 \label{eq:bridge-reuse}
\end{equation}
Historical mass coefficients are multiplied by \(c\), while a
flux package already containing \(h_{\rm eff}\) must not receive
that factor a second time. Previously precontracted nonlinear
histories depend on the old inverse matrix and must be remapped from
their original flux spectra at the interface. Reusing those old
contracted vectors directly would change the algorithm.

Exact-arithmetic reuse does not imply bitwise equality across
implementations. The archive verifies byte-for-byte equality of 936
arrays in 24 fixed comparisons, retaining the original floating-point
expression grouping and environment. Changing the reduction order,
library, spatial basis, lifting, or adaptation policy requires a new
check. Shared preparation removes repeated spatial work but leaves
the phase-specific factorizations, operator packages, transfers, and
execution graphs; it is not by itself a timing result.

\section{Common phase coordinates and the limits of algebraic improvement}
\label{sec:phase-analysis}

\subsection{Simultaneous congruence for one angular mode}

For one fixed mode, let the phase matrices be
\(A_j=a_jM+h_jK\), with \(a_j,h_j>0\). The two constant segments
and the bridge have this form. We use a different symbol for the
reference matrix, \(H_*\), to distinguish it from the early time
step \(H\). Assume that \(M,K\) are real symmetric positive
semidefinite and
\(H_*=a_*M+h_*K\) is positive definite for \(a_*,h_*>0\).
Equivalently, \(M\) and \(K\) have no common nonzero null vector.
The standard symmetric-definite generalized eigenproblem
\cite{r10} provides the following representation.

\paragraph{Proposition 3 (common coordinates).}
There exists a real invertible \(P\) such that
\(P^{\mathsf T}H_*P=I\),
\(P^{\mathsf T}MP=\operatorname{diag}(\lambda_i)\), and, with
\(\theta_i=a_*\lambda_i\),
\begin{equation}
 \begin{gathered}
 0\leq\theta_i\leq1,\qquad
 P^{\mathsf T}A_jP=\operatorname{diag}(d_{ji}),\\
 d_{ji}=\frac{a_j}{a_*}\theta_i+
             \frac{h_j}{h_*}(1-\theta_i)>0.
 \end{gathered}
 \label{eq:phase}
\end{equation}

\paragraph{Proof.}
Write \(H_*=LL^{\mathsf T}\). The symmetric positive
semidefinite matrix \(L^{-1}ML^{-\mathsf T}\) has an orthogonal
eigenbasis \(Q\). For \(P=L^{-\mathsf T}Q\), its eigenvalues
give \(P^{\mathsf T}MP=\operatorname{diag}(\lambda_i)\) and
\(P^{\mathsf T}H_*P=I\). Further,
\(h_*P^{\mathsf T}KP=I-a_*\operatorname{diag}(\lambda_i)\)
is positive semidefinite, so \(0\leq a_*\lambda_i\leq1\).
Substitute the linear combination defining \(A_j\) to obtain
\eqref{eq:phase}. Every diagonal value is a convex combination of
two positive numbers. \hfill\(\square\)

Consequently,
\begin{equation}
 A_j^{-1}b=P\operatorname{diag}(d_{ji}^{-1})P^{\mathsf T}b.
 \label{eq:phase-inverse}
\end{equation}
This is a congruence transformation, rather than a claim that the
original-coordinate matrices commute. The construction also applies
to the separate mean matrices under the same hypotheses. Although
an Euler matrix has positive coefficients and fits the algebra, the
archived common-phase prototype retains the original Euler startup;
that phase was not independently promoted by the prototype tests.

\subsection{A geometric anchor for a worst-case spectral bound}

Set \(r_j=h_j/a_j\), \(\eta=h_*/a_*\),
\(r_{\min}=\min_jr_j\), and \(r_{\max}=\max_jr_j\).
After extracting \(a_j/a_*\) from \(d_{ji}\), its remaining
factor is \(\theta_i+(r_j/\eta)(1-\theta_i)\), which lies
between \(1\) and \(r_j/\eta\). Therefore
\begin{equation}
 \kappa_2(P^{\mathsf T}A_jP)
   \leq B_j(\eta)=\max\{r_j/\eta,\eta/r_j\}.
 \label{eq:bound}
\end{equation}
For a finite set of positive phases,
\begin{equation}
 \eta_* = \sqrt{r_{\min}r_{\max}},\qquad
 \min_{\eta>0}\max_j B_j(\eta)
       =\sqrt{r_{\max}/r_{\min}}.
 \label{eq:anchor}
\end{equation}
Indeed, the worst bound is
\(\max\{r_{\max}/\eta,\eta/r_{\min}\}\). The two positive
terms have fixed product \(r_{\max}/r_{\min}\); their maximum is
smallest when they are equal, proving \eqref{eq:anchor}.

The distinction between the bound and the actual spectrum is
essential. If the \(\theta_i\) do not approach both endpoints, the
interval estimate is loose. The anchor minimizes that estimate, not
necessarily the actual maximum condition number for the given
matrices. Rescaling \(a_*,h_*\) together does not change this bound
but does change diagonal magnitudes, so the rule does not imply
that all denominators are close to one. It also does not improve
\(\kappa_2(A_j)\) in the original coordinates or bound the errors
of applying \(P\) and \(P^{\mathsf T}\).

The geometric-anchor statement is an algebraic derivation, separately
reviewed, and was not a timed candidate. The measured setup overhead
of the earlier common-phase prototype cannot be assigned to this
untested choice of anchor. Neither the derivation nor the prototype
constitutes a random-feature approximation theorem or a stability
theorem for the time integrator.

\subsection{Why common coordinates need not yield a faster solver}

Across angular modes, \(M_k=M^{(0)}+k^2M^{(1)}\) and
\(K_k=K^{(0)}+k^2K^{(1)}+k^4K^{(2)}\). Polynomial dependence on
\(k^2\) does not imply a common eigenbasis. In a fixed whitening
metric, real symmetric coefficient matrices admit a common
orthogonal diagonalization precisely when they commute pairwise.
The archived normalized commutators, about \(0.051\)--\(0.076\),
do not support replacing all modes by one exact common transform.

Even within one mode, the storage contract matters. Suppose the
device expects precontracted dense operators \(A_j^{-1}B\), where
\(B\) contains mass and nonlinear projection columns. Keeping this
interface requires
\(P\operatorname{diag}(d_j^{-1})P^{\mathsf T}B\) for every phase.
For rank \(d\), projection width proportional to \(q\), and
\(p_{\rm ph}\) phases, this entails an \(O(d^3)\)
eigendecomposition and approximately \(p_{\rm ph}+1\)
\(O(d^2q)\) dense multiplications. Direct phase preparation uses
\(p_{\rm ph}\) factorizations and phase-wise applications instead.
The common representation still writes all phase packages. Operation
counts, library throughput, conditioning, and memory layout therefore
have to be evaluated together.

A different design could keep the time-dependent coefficients in
the \(P\) coordinates and use diagonal phase updates throughout.
It would also require mode-dependent transformed reconstruction,
lifting, nonlinear projection, and observation operators. The current
trajectory evidence does not validate that design. It remains a
specific engineering opportunity, rather than an inferred speedup
from \eqref{eq:phase-inverse}.

\subsection{Separate approximation, integration, and implementation errors}

Four levels are relevant: the continuous finite-domain solution; a
quadrature-defined spatial semidiscretization; its exact-arithmetic
time-discrete solution; and the executed floating-point solution.
They distinguish domain error, trial-space and quadrature error,
time-discretization error, and implementation error. On a common
physical reconstruction grid, the triangle inequality gives an
additive decomposition in a common absolute norm. Relative error
bounds require explicit denominator conversions; two percentages
with different reference fields cannot simply be added.

The observed field gates in this paper compare numerical trajectories
to two resolved numerical references. The implementation gate compares
GPU and CPU realizations of the same discrete method. Neither is a
rigorous upper bound on error against the unknown continuous
solution. In addition, a small velocity error does not control its
spatial derivatives: \eqref{eq:kinematics} differentiates the
streamfunction twice for vorticity. Vorticity and near-field gates
must therefore be imposed independently.

For completeness, a conditional perturbation estimate explains why
small operator residuals still require trajectory checks. Consider
two executions of the same step and lifting, with coefficient errors
\(e_n\) and \(e_{n-1}\), and an additive algebraic residual
\(r_{n+1}\) in the step-scaled linear equation. If the nonlinear
map is Lipschitz with constant \(L_f\) on a region containing both
trajectories, then a compatible coefficient norm gives
\begin{equation}
\begin{aligned}
 \|e_{n+1}\|\leq\|A_n^{-1}\|\big[&
 (b_0\|M\|+h(1+\rho)L_f)\|e_n\|\\
 &+(b_1\|M\|+h\rho L_f)\|e_{n-1}\|
 +\|r_{n+1}\|\big].
\end{aligned}
 \label{eq:perturbation-accounting}
\end{equation}
Subtracting the two updates and applying the norm and local Lipschitz
inequality proves this estimate. It is only error bookkeeping for
the same finite-dimensional scheme; its constants may be very
pessimistic, and it does not establish a uniform long-time bound.
It explains why agreement of an isolated FFT, small linear residuals,
or a favorable transformed condition number cannot replace the
physical accuracy checks at the full integration horizon.
